\documentclass[11pt]{article}
\usepackage[T1]{fontenc}
\usepackage{lmodern,microtype}
\usepackage[margin=0.91in]{geometry}
\usepackage{amsmath,amssymb,amsthm,mathtools}
\usepackage{booktabs,longtable,array,enumitem}
\usepackage{xcolor,hyperref,fancyhdr,listings}
\definecolor{ink}{HTML}{173B55}
\definecolor{muted}{HTML}{52616B}
\hypersetup{colorlinks=true,linkcolor=ink,urlcolor=ink,citecolor=ink,pdftitle={Sparse support packing for Diophantine consistency equations},pdfauthor={AI-assisted research prepared for Vladimir Reshetnikov}}
\pagestyle{fancy}\fancyhf{}\lhead{\small Sparse support and ZFC consistency}\rhead{\small 23 October 2026}\cfoot{\thepage}
\setlength{\headheight}{14pt}
\setlength{\parskip}{4pt}\setlength{\parindent}{0pt}
\setlist{itemsep=2pt,topsep=4pt}
\newtheorem{theorem}{Theorem}[section]
\newtheorem{lemma}[theorem]{Lemma}
\newtheorem{proposition}[theorem]{Proposition}
\newtheorem{corollary}[theorem]{Corollary}
\theoremstyle{definition}\newtheorem{definition}[theorem]{Definition}
\newcommand{\Z}{\mathbb Z}\newcommand{\N}{\mathbb N}\newcommand{\Pos}{\mathbb Z_{>0}}
\newcommand{\Con}{\operatorname{Con}(\mathrm{ZFC})}
\newcommand{\ZRef}{\operatorname{Ref}_{\mathrm{ZFC}}}
\newcommand{\wt}{\operatorname{wt}}
\newcommand{\pop}{\operatorname{pop}}
\newcommand{\coeff}{\operatorname{coeff}}
\newcommand{\sm}{\mathsf}
\newcommand{\pin}{\texttt{124c07c0b77a9e20e84676514d8b1ff78536ea3f}}
\newcommand{\repo}{https://github.com/VladimirReshetnikov/ProveIt}
\newcommand{\base}{https://github.com/VladimirReshetnikov/ProveIt/blob/124c07c0b77a9e20e84676514d8b1ff78536ea3f/Computability/DiophantineConsistency}
\lstset{basicstyle=\ttfamily\footnotesize,breaklines=true,columns=fullflexible,frame=single,rulecolor=\color{muted},showstringspaces=false}
\begin{document}
\begin{center}
{\color{ink}\LARGE\bfseries Sparse support packing for explicit\\[3pt] Diophantine consistency equations}\par
\vspace{9pt}
{\large Nine positive unknowns; ten unrestricted integer unknowns}\par
\vspace{6pt}
{Research prepared for Vladimir Reshetnikov and the ProveIt project}\par
{23 October 2026}
\end{center}

\begin{abstract}
The latest multi-target construction in ProveIt distinguishes all 528
monomials of degree at most two in 31 variables, although its 22 source
equations use only 71 distinct monomials. We replace complete quadratic
separation by the weaker, sufficient requirement that each of these 71
monomials be isolated from every possible contaminating monomial. Explicit
certificates reduce the largest weight from 799 to 239 or 220. A second
change replaces a cubic variable base by a fixed multiple of a square.
Together these changes produce a nine-positive-variable polynomial of
degree 23,878,392 with 670 live binary arithmetic operations, and a
651-operation alternative of degree 25,936,016. The compared construction
has degree 129,875,980 and costs 788 operations. Applying the sparse weights
to the existing unrestricted-integer reduction gives ten-variable
polynomials of degree 65,506,504 in 675 operations, or degree 71,151,176 in
662 operations, compared with 237,520,456 in 782 operations.

The new contribution is an explicit, audited simplification of the source
encoding. The proof transfers solvability to the same fixed quadratic
system used by the repository. Its interpretation as ZFC refutability
uses the inherited ordinary mathematical machine/compiler theorem; the
ten-integer route also uses Sun's cited reduction. No new end-to-end
Lean or Rocq certificate, proof of consistency, or global minimality
theorem is claimed. The archive contains the exact finite coefficient
payloads, executable polynomial circuits, proofs, and independent checks.
\end{abstract}

\section{Result, meaning, and comparison}

Write $\ZRef$ for the arithmetical assertion that there is a finite formal
ZFC proof of a fixed contradiction, and put $\Con:=\neg\ZRef$. We seek a
\emph{fixed} polynomial, not a family with an unspecified index or a
coefficient chosen according to whether ZFC is consistent.
All coefficients in this construction are supplied by finite recipes and
materialized binary integers. No recipe tests consistency.

The source snapshot is \pin. It includes the wave-6 packages 38--45,
although the collection README at that snapshot still mainly describes
wave 5. The comparison below uses the actual wave-6 files
\cite{source40,source45}, rather than treating that README as current.

\begin{center}\small
\begin{tabular}{@{}llrr@{}}
\toprule
Domain & Construction & Total degree & Binary operations\\
\midrule
$\Pos^9$ & Package 45 baseline & 129,875,980 & 788\\
$\Pos^9$ & New, weights $W=239$ & 25,936,016 & 651\\
$\Pos^9$ & New, weights $W=220$ & 23,878,392 & 670\\
\midrule
$\Z^{10}$ & Package 40 baseline & 237,520,456 & 782\\
$\Z^{10}$ & New, weights $W=239$ & 71,151,176 & 662\\
$\Z^{10}$ & New, weights $W=220$ & 65,506,504 & 675\\
\bottomrule
\end{tabular}
\end{center}

Both new weight choices improve both displayed measures over their
respective baselines. Within each new pair, the degree and operation
objectives differ: $W=220$ has lower degree, while $W=239$ is cheaper to
evaluate. The latter has a more favorable addition-chain and expression
sharing structure; a smaller exponent does not automatically imply fewer
gates.

For nine positive variables, the degree improvement is a factor of
approximately $5.44$ at $W=220$, and the operation reduction is approximately
$17.4\%$ at $W=239$. For ten unrestricted integers, the corresponding
figures are approximately $3.63$ and $15.3\%$.
These are substantial reductions in the low-variable part of the
repository's frontier. They do not beat the separate 67-operation
higher-variable equation or the quartic alternatives on their preferred
measures.

\paragraph{Arithmetic convention.}
Every live binary $+$, $-$, or $\times$ costs one operation, including
multiplication by a fixed numeral. Reusing a register is free. Fixed
powers are emitted as multiplication gates. A displayed abbreviation is
a polynomial expression, not a new existential variable. Degree is total
degree in the quantified variables; fixed integer coefficients have degree
zero. The JSON circuits are complete division-free descriptions of ordinary
polynomials in $\Z[X_1,\ldots,X_n]$.

\paragraph{Exact scope of the theorem.}
The unconditional new transfer result is
\[
  \exists\boldsymbol{x}\in\Pos^9\ F_{9,W}(\boldsymbol{x})=0
  \quad\Longleftrightarrow\quad
  \exists\boldsymbol{u}\in\Pos^{31}\
       \bigwedge_{j=1}^{22}r_j(\boldsymbol{u})=0.
  \tag{1}\label{eq:transfer}
\]
The right side is the \emph{same fixed system}, with the same complete
four-coefficient tuple, as package 31. Composing (1) with its inherited
machine/compiler interpretation gives the asserted consistency
equivalence, by ordinary mathematics formalizable in ZFC:
\[
 \mathrm{ZFC}\ \vdash\
 \left(\neg\exists\boldsymbol{x}\in\Pos^9\ F_{9,W}(\boldsymbol{x})=0\right)
 \longleftrightarrow\Con.
 \tag{2}\label{eq:zfc}
\]
The analogous composition applies to the ten-integer construction.
Equation (2) inherits the ordinary compiler proof as a mathematical
dependency; this note does not turn the repository's open formal
compiler and object-language proof-export obligations into completed
kernel certificates. Section~\ref{sec:boundary} records that distinction
explicitly.

\section{The fixed quadratic source}
\label{sec:parent}

The 31 positive source coordinates, in the order used by the weight
vectors, are
\[
\begin{gathered}
 J,F,\alpha,z,f,h,i,o,j,s,w,\tau,\eta,\zeta,y,\gamma,\\
 N,E,P,c,a,X,D,C_2,L,F_2,V,g,I,B_{\rm aux},\Delta.
\end{gathered}
\]
The four fixed ports are $(\beta,K_0,M_C,M_F)$ from package 31's
\texttt{wang943-mixed-121x8801} instance. The native base is
$\mathcal B=2^{1,064,921}$, and $\beta=\mathcal B-1$.
Here $M_F$ is the \emph{shifted} high mask. Mixing this port with an
unshifted mask, or taking ports from different compiler tuples, would
define a different problem.

Define only for writing the source equations
\[
\begin{gathered}
q=\beta J+1,\quad k=\eta+\zeta,\quad H=4a+3,\quad C=q-F-\alpha,\\
U=L+1,\quad Y=a-E,\quad Q=c-\eta,\quad T=P-k,\quad R=jc-V.
\end{gathered}
\]
The ordered equations are
\begin{alignat*}{3}
 U&=q^2, &\qquad N&=qU, &\qquad X&=wq,\\
 Y&=sN, & E&=XY, & Q&=kY,\\
 T&=EQ, & \tau^2&=PT+1, & D&=X+ac+\gamma H,\\
 \Delta&=a^2+H, & C_2&=c^2, & D^2&=\Delta C_2+1,\\
 k&=R+1+hE,&&&&
\end{alignat*}
\begin{align*}
R&=(q-F)(N-q)-CL+(M_C+qM_F)J,\\
(K_0+w)C&=F+z\beta J,\\
F_2&=f^2,\qquad V=of-c,\qquad g=y^2-V^2,\\
I&=iC_2,\qquad B_{\rm aux}=I^2,\qquad
B_{\rm aux}=\Delta(F_2-1),\\
y^2-B_{\rm aux}g&=1.
\end{align*}
All 22 left-minus-right residuals are quadratic. This is the exact
60-operation comparison system in
\texttt{quartic31\_system60.json}, whose SHA-256 is
\begin{center}\footnotesize\ttfamily
ac2a67782d344d1e74fd93bc3f09fae0097eaed1829fd8194ae235ae84763e48.
\end{center}
The archive preserves that source file and the complete coefficient
payloads. The port bit lengths are respectively
\[
 1,064,921,\quad1,070,004,\quad1,064,921,\quad1,064,922.
\]
Rebuilding the tuple from the literal machine transitions and 943 tiles
reproduces the pinned hashes. Its four uncompressed payloads total
533,099 bytes; their compressed forms total 10,485 bytes.
These are the native inputs, not the coefficient list of a fully
expanded final polynomial.

Set each positive source coordinate to $x_i+1$ and regard $x_i\geq0$.
Reorder the rows as
\[
 1,2,3,4,6,7,\ldots,22,5.
\]
Call the reordered shifted residuals $p_1,\ldots,p_{22}$.
The last row has constant coefficient 1. This has two uses: the tuple
$\boldsymbol{x}=0$ is not a common zero, and the top coefficient of
the encoding polynomial below is positive and independent of the ports.
Independent sparse polynomial replay gives 136 row terms, with 71
distinct monomials in their union: one constant, 31 linear monomials,
and 39 quadratic monomials.

\section{Only used monomials need to be separated}
\label{sec:sparse}

Let
\[
 \mathcal A=\bigcup_j\operatorname{supp}(p_j),\qquad
 \mathcal Q=\{\alpha\in\N^{31}:|\alpha|\leq2\}.
\]
Thus $|\mathcal A|=71$ and $|\mathcal Q|=\binom{33}{2}=528$.
The condition actually needed is
\[
 \alpha\in\mathcal A,\ \delta\in\mathcal Q,\quad
 \alpha\cdot\boldsymbol w=\delta\cdot\boldsymbol w
 \quad\Longrightarrow\quad\alpha=\delta.
 \tag{3}\label{eq:isolation}
\]
There is no requirement that two monomials outside $\mathcal A$ have
different weights. Their collisions never select a nonzero source
coefficient. Full $B_2$ separation is sufficient but stronger than (3).

Here are the two delivered vectors, in the coordinate order of
Section~\ref{sec:parent}:
\begin{align*}
\boldsymbol w_{220}={}&(24,89,95,146,205,155,147,197,203,78,29,\\
 &148,47,208,25,217,129,169,30,215,220,219,135,\\
 &166,27,133,32,200,43,204,152),\\[3pt]
\boldsymbol w_{239}={}&(237,231,223,221,233,193,195,215,120,190,226,\\
 &142,125,151,148,169,214,134,140,121,126,172,128,\\
 &196,238,191,239,192,197,184,205).
\end{align*}
An exact finite certificate checks (3) against all 528 expansion
monomials for every used monomial. The optimizer is not trusted: any
method may propose the vector, and a short independent verifier checks
the property. No minimality of either vector is asserted.

For either vector put
\[
 W=\max_iw_i,\quad d=2W+1,\quad h_j=jd-1,\quad
 H_* =22d-1,\quad K_*=H_*+2W.
\]
If $p_j(\boldsymbol x)=\sum_\alpha p_{j,\alpha}\boldsymbol x^\alpha$,
define the fixed integers
\[
 c_{j,\alpha}=(2-|\alpha|)!\left(\prod_i\alpha_i!\right)p_{j,\alpha}
\]
and the univariate polynomial
\[
 \mathcal D(Z)=\sum_{j,\alpha}c_{j,\alpha}
                       Z^{h_j-\alpha\cdot\boldsymbol w}.
 \tag{4}\label{eq:D}
\]
All exponents in (4) are nonnegative fixed integers. The four-port
expressions for the 136 coefficients are given exactly in the data files.
The first target is $2W$, the lowest position that uniformly accommodates
every quadratic term. This one-place shift saves three operations in
each delivered nine-variable evaluation schedule.

\begin{lemma}[Sparse extraction]
For every integer tuple $\boldsymbol{x}$ and every $j$,
\[
 [Z^{h_j}]\ \mathcal D(Z)
       \left(1+\sum_i x_iZ^{w_i}\right)^2=2p_j(\boldsymbol{x}).
 \tag{5}\label{eq:extract}
\]
\end{lemma}
\begin{proof}
A source term in row $k$ and an expansion monomial $\delta$ reach
$h_j$ exactly when
\[
 (k-j)d=\alpha\cdot\boldsymbol w-\delta\cdot\boldsymbol w.
\]
The right side has absolute value at most $2W$, so $k=j$.
Property (3) then forces $\alpha=\delta$. The coefficient in the
squared expansion is
$2!/((2-|\alpha|)!\prod_i\alpha_i!)$; multiplying by the coefficient
in (4) gives exactly $2p_{j,\alpha}$.
\end{proof}

The intervals $[h_j-2W,h_j]$ for the different source rows are disjoint.
Within each row, (3) makes its used exponents distinct. Therefore every
coefficient of $\mathcal D$ is one $c_{j,\alpha}$, not a hidden sum of
several source coefficients. Its highest coefficient, at $H_*$, is 2.
Moreover, if $|c_{j,\alpha}|<M$, every coefficient of the product in
(5) has absolute value at most
\[
 M\left(1+\sum_i|x_i|\right)^2.                 \tag{6}\label{eq:coefbound}
\]
This bound is unchanged by collisions among unused expansion monomials:
one sums their multinomial coefficients, not the number of distinct
weights.

\section{A fixed multiple of a square replaces the cubic base}
\label{sec:base}

The supplied coefficient recipes and the port bit lengths give the
explicit bound
\[
 M=2^{2,134,927},\qquad \lambda=4096M=2^{2,134,939}.
 \tag{7}\label{eq:lambda}
\]
The verifier computes the bound from the integer recipes, rather than
accepting the exponent as an unexplained parameter. Work temporarily
with nine coordinates
\[
 \varepsilon,\phi,n\geq0,\qquad g,h,s,w,i,j\geq1.
 \tag{8}\label{eq:mixed}
\]
Set
\[
 b=\varepsilon+2,\qquad B=\lambda b^2.
 \tag{9}\label{eq:base}
\]
In the old construction $b$ had a large additive floor and $B=b^3$.
The reason for that cubic was a coefficient bound proportional to
$Mb^2$. Since $M$ is fixed, placing it in a fixed multiplicative
coefficient suffices: $1024Mb^2=B/4$. This removes one unit of
variable degree from every packed power without introducing a variable.
The fixed value $\lambda$ is exactly package 45's former additive
\texttt{base\_floor}, now used as a multiplicative \texttt{base\_scale}.
The other 16 fixed leaves are unchanged. Thus this improvement does not
introduce a larger fixed coefficient or enlarge the fixed-leaf payload:
the 17 compressed nine-variable payloads total 64,286 bytes.

Define the following polynomial registers:
\begin{align*}
 \ell&=\sum_{r=1}^{31}B^{w_r}, & D_c&=\mathcal D(B),\\
 N_0&=B^{W+1}, &N_1&=4B^{K_*+1},\qquad N=N_0N_1,\\
 M_c&=N_0-1-(b-1)\ell, &T_*&=\sum_{r=1}^{22}B^{h_r},\\
 T_c&=M_c+(B-2)T_*N_0,\\
 J_c&=2(1+g)^2D_c+B\sum_{r=0}^{K_*}B^r,\\
 S_c&=g+J_cN_0,\\
 R&=(N-1)\bigl((S_c+T_c+1)N+T_c+1\bigr),\\
 G_c&=32bB^W.
\end{align*}
The subscript $c$ distinguishes coding registers from similarly named
quantities in the later Pell construction. Every sum is finite with a
fixed upper endpoint. Geometric sums are computed polynomially; a
fraction such as $(B^{k+1}-1)/(B-1)$ is not an instruction in the circuit.

\begin{lemma}[Uniform outer bounds]
On every assignment (8), $D_c,M_c,T_c,J_c,S_c>0$, and
$R,N\geq8$, $b\leq R,N$. If also $g<4bB^W$, then
\[
 0<g<N_0,\qquad0<J_c<N_1,\qquad0<S_c,T_c<N.
 \tag{10}\label{eq:outer}
\]
\end{lemma}
\begin{proof}
Since $B-1>M$ and the top coefficient of $\mathcal D$ is 2,
\[
 D_c>2B^{H_*}-M\sum_{r<H_*}B^r>B^{H_*}>0,
 \qquad D_c<\frac{MB^{H_*+1}}{B-1}.
\]
Also $\ell<(B^{W+1}-1)/(B-1)$ and $b-1<B-1$, so $M_c>0$.
The other unconditional assertions follow immediately from the displayed
positive formulas. On the bounded-$g$ branch, $4b<B$ and
$1+g\leq4bB^W$, so
\[
 2(1+g)^2D_c<\frac{32Mb^2}{B-1}B^{K_*+1}.
\]
Here $32Mb^2=B/128$. Adding the geometric-mask term yields
$0<J_c<4B^{K_*+1}=N_1$. The target-mask bound
$(B-2)T_*<B^{H_*+1}$ gives $T_c<N$, while $g<N_0$ and
$J_c<N_1$ give $S_c<N$.
\end{proof}

\section{Why the carry code imposes every source equation}
\label{sec:carry}

Assume first that $b$ is a power of two and $g<4bB^W$.
Then $B,N_0,N_1,N$ are powers of two. For $N=2^r$ and
$0\leq S,T<N$, the identity
\[
 (N-1)((S+T+1)N+T+1)
 =(S+T)N^2+(N-1-S)N+(N-1-T)
 \tag{11}
\]
has nonoverlapping binary blocks. Thus
\[
 \pop(R)=2r+\pop(S_c+T_c)-\pop(S_c)-\pop(T_c).
\]
Legendre's factorial identity gives
$v_2\binom{2R}{R}=\pop(R)$. Consequently
\[
 N^2\mid\binom{2R}{R}\quad\Longleftrightarrow\quad
 \text{adding $S_c$ and $T_c$ produces no binary carry}.
 \tag{12}
\]

The low-block condition is that $g$ and $M_c$ have disjoint 1-bits.
The base-$B$ digit of $M_c$ is $B-b$ at each allowed $w_i$ and $B-1$
at every other position through $W$. Therefore this condition is
equivalent to
\[
 g=\sum_i x_iB^{w_i},\qquad0\leq x_i<b.
 \tag{13}
\]
The bound $g<N_0$ excludes any additional higher digit.

For such $g$, let $u_r$ be the coefficients of
$\mathcal D(Z)(1+\sum_i x_iZ^{w_i})^2$. Its degree is at most $K_*$.
By (6),
\[
 |u_r|\leq M(32b)^2=B/4<B/2.
\]
Hence $a_r=u_r+B/2$ satisfies $0<a_r<B$ and
\[
 J_c=2\sum_{r=0}^{K_*}a_rB^r,
 \qquad (B-2)T_*=2\sum_{j=1}^{22}(B/2-1)B^{h_j}.
\]
Dividing both operands by 2 is a common binary shift. The target digit
$B/2-1$ uses all the low bits of its cell. Disjointness forces the positive
digit $a_{h_j}$ to be exactly $B/2$. By (5), this is precisely
$p_j(\boldsymbol x)=0$, separately for each of the 22 rows.
There is no inference from a sum of residuals to their separate vanishing.

Conversely, given a common nonnegative zero $\boldsymbol x$ of the source,
choose any sufficiently large power of two $b>\max_i x_i$, with $b\geq2$.
Equation (9) supplies a power-of-two $B$. Set
$\varepsilon=b-2$ and define $g$ by (13). Since the all-zero tuple is not
a source zero, $g>0$. Moreover $g<(b-1)\sum_{r\leq W}B^r<bB^W$.
Both no-carry conditions hold, and therefore
\[
 N^2\mid\binom{2R}{R},\quad b\text{ is a power of two},\quad
 32g<G_c.                                             \tag{14}
\]
The weaker inequality $8g<G_c$ supplied by the soundness direction of
the ratio theorem gives exactly the bounded-$g$ hypothesis used above.
This ordering avoids circular use of the digit bound.

\section{The nine-variable polynomial, explicitly}
\label{sec:equation}

Use the coding registers just defined, and the remaining coordinates
$h,s,w,i,j>0$, $\phi,n\geq0$. Define
\begin{align*}
U&=N^2w,&Y&=N^2s,&M_0&=RY,\\
P&=2M_0^2U,&K&=R+1+h(P-1),&A&=M_0(U+1),\\
B_1&=2R+1,&W_1&=bw,&C&=B_1+W_1+\phi,\\
J&=4A-5,&Q&=3W_1C-2(W_1^2-1),\\
D&=(A^2-1)C^2+1,&E&=iJDC^2,&F&=(A^2-1)E^2+1,\\
G&=A+F(CD+1-A),&H&=B_1+jC,&I&=(G^2-1)H^2+1,\\
a_0&=(P^2-1)K^2+1,&a_1&=DFI,\\
u&=JF,&v&=QF+J(H-C),\\
m&=K^2G_c-4(C-KY)^2G_c-8gK^2.
\end{align*}
The numerical Jones ratio theorem states that, under the unconditional
bounds already proved, the four conditions
\[
 a_0,a_1\text{ are integer squares},\quad u\mid v,\quad m>0
 \tag{15}\label{eq:ratio}
\]
imply the central-binomial divisibility, the power-of-two condition on
$b$, and $8g<G_c$. Conversely (14) supplies witnesses satisfying (15).
The precise numerical theorem is
\texttt{ratio\_sound} and \texttt{ratio\_complete} in namespace
\texttt{Jones1982.Nine} of the pinned \texttt{NineRatio.lean};
it does not fix the old positional weights
\cite{jones,ratio}. Thus changing the source coding does not silently
change a hypothesis of the ratio theorem.

To express (15) with the one remaining coordinate $n$, put
\[
 V_0=PK,\qquad V_1=(AC)(AE)(GH),\qquad M_d=Q+H,
\]
\[
 \Omega=um,\qquad
 \Theta=u(n+M_d)+v-\Omega(2M_d+V_0V_1).
\]
The equation is
\[
 \boxed{\ F_{9,W}=
 \bigl[\Theta^2-\Omega^2(a_0+V_0^2a_1)\bigr]^2
 -4\Omega^4V_0^2a_0a_1=0.\ }                    \tag{16}\label{eq:F9}
\]
Together with Sections~\ref{sec:parent}--\ref{sec:base}, this is a
finite definition of the polynomial. It contains no variable exponents,
divisions, binomial coefficients, square-root functions, or extra
quantifiers. Those objects occur only in its correctness proof.

\begin{lemma}[Finalizer]
For the preceding outer parameters, an $n\geq0$ satisfying (16) exists
if and only if (15) holds.
\end{lemma}
\begin{proof}
All of $P,K,A,C,J,D,E,F,G,H,I$ are positive, and the formulas give
\[
0<a_0<V_0^2,\qquad0<a_1<V_1^2,\qquad
0<v/u<M_d-1.                                      \tag{17}
\]
For example $V_0^2-a_0=K^2-1>0$; the three inequalities
$(AC)^2>D$, $(AE)^2>F$, and $(GH)^2>I$ give the second bound.
Also $H-C=B_1+(j-1)C>0$, and
$v/u=Q/J+(H-C)/F<Q+H-1$.

If $m=0$, then $\Omega=0$ and
$\Theta/u=n+M_d+v/u>0$, so (16) cannot hold. If $m\ne0$,
the four-factor identity for (16) gives
\[
 \Theta/\Omega=r_0+V_0r_1,\qquad r_0^2=a_0,\quad r_1^2=a_1.
\]
The roots are real. The rational number on the left cannot vanish:
otherwise $|r_1|=|r_0|/V_0<1$, contradicting $a_1\geq1$.
Writing it as $x\ne0$ gives
$2xr_0=x^2+a_0-V_0^2a_1$, so $r_0$ and then $r_1$ are rational.
A rational square root of an integer is integral.

By (17), $|r_0|\leq V_0-1$ and $|r_1|\leq V_1-1$. Hence
\[
 Z=2M_d+V_0V_1+r_0+V_0r_1\geq2M_d+1.
\]
The vanishing factor says
$n+M_d+v/u=mZ$. Its left side is positive, so $m>0$.
Its right side is an integer, so $u\mid v$.
Conversely, if (15) holds, choose integral roots and set
$n=mZ-M_d-v/u$. The same bounds make this an integer at least 2,
and the corresponding factor vanishes.
\end{proof}

Composition with the carry-code lemma proves (1). To obtain exactly nine
\emph{positive} input variables, replace $\varepsilon,\phi,n$ by their
positive input coordinates minus 1. To obtain nine \emph{nonnegative}
input variables instead, shift $g,h,s,w,i,j$ upward by 1.
These are different supplied circuits. The positive circuit must not be
interpreted as an unrestricted-integer equation.

\section{Exact degree and why the leading cancellation matters}
\label{sec:degree}

Let
\[
 \zeta=H_*+3W+2=47W+23,\qquad t=\deg N=2\zeta=94W+46.
\]
On the diagonal where each quantified variable equals an indeterminate
$T$, the leading terms of $N$ and $R$ are
$4\lambda^{\zeta}T^t$ and
$16\lambda^{3\zeta}T^{3t}$.
The high geometric term in $J_c$ dominates its sparse-coefficient term,
and the $S_cN^2$ term dominates $R$.

The following table records the degrees needed for an analytic
certificate. All are also full multivariate upper bounds.
\begin{center}\small
\begin{tabular}{@{}lr@{\qquad}lr@{}}
\toprule
Register & Degree & Register & Degree\\
\midrule
$P$ & $12t+3$ & $K$ & $12t+4$\\
$A$ & $7t+2$ & $C$ & $3t$\\
$D$ & $20t+4$ & $E$ & $33t+7$\\
$F$ & $80t+18$ & $G$ & $103t+22$\\
$H$ & $3t+1$ & $I$ & $212t+46$\\
$a_0$ & $48t+14$ & $a_1$ & $312t+68$\\
$u$ & $87t+20$ & $m$ & $28t+2W+11$\\
$\Omega$ & $115t+2W+31$ & $V_0$ & $24t+7$\\
$V_1$ & $156t+34$ & $M_d$ & $3t+2$\\
\bottomrule
\end{tabular}
\end{center}

Naive degree propagation through (16) overestimates the result, because
the highest terms of $V_1^2$ and $a_1$ cancel. Resolve this by the exact
identity
\begin{align*}
\Delta_V:=V_1^2-a_1={}&(C^2-1)(AE)^2(GH)^2\\
 &+D(E^2-1)(GH)^2+DF(H^2-1).
\end{align*}
Its first two summands have degree $298t+64$ and each has leading
coefficient $2^{518}\lambda^{298\zeta}$; their sum does not cancel.
With $E_0=u(n+M_d)+v-2\Omega M_d$, the inner bracket in (16) equals
\[
 E_0^2-2E_0\Omega V_0V_1+\Omega^2V_0^2\Delta_V-\Omega^2a_0.
\]
The third term uniquely dominates. It has degree
$576t+4W+140$ and leading coefficient
$2^{1025}\lambda^{576\zeta+2W}$.
The final term subtracted in (16) has strictly lower degree.
Therefore
\[
 \boxed{\deg F_{9,W}=1152t+8W+280=108296W+53272.}       \tag{18}
\]
The exact nonzero diagonal leading coefficient is
\[
 2^{2050}\lambda^{1152\zeta+4W}
 =2^{2050}\lambda^{54148W+26496}.                      \tag{19}
\]
Thus the upper bound is attained for the actual fixed coefficients, not
merely for a generic formal coefficient tuple. The separately exported
cancellation-resolved circuit checks the same identity. Modular leading
term checks provide additional independent transcription evidence; (19)
is the mathematical reason the leading term is nonzero.

The effects of the changes can be separated:
\begin{center}\small
\begin{tabular}{@{}lrr@{}}
\toprule
Source encoding & Maximum weight & Nine-variable degree\\
\midrule
Baseline, cubic base &799&129,875,980\\
Only change to scaled square base &799&86,584,080\\
Only use sparse weights &220&35,820,904\\
Sparse weights and scaled square base &220&23,880,696\\
Also shift all targets down by one &220&23,878,392\\
\bottomrule
\end{tabular}
\end{center}
The coefficient heights of expanded polynomials are enormous in all
these representations. We compare exact circuits and fixed coefficient
payloads, not an ungenerated dense monomial list.

\section{Ten unrestricted integer variables}
\label{sec:ten}

The same sparse support certificate can be inserted into package 40's
quadratic source and its existing ten-integer tail. That baseline already
uses a fixed multiple of a square as its packing base, and already uses
the quintic divisibility merger
\[
 (r^2+23S^2)(r^3-rS^2-S^3),\qquad r=T+SY.
\]
Neither device is a new contribution of this note on the ten-integer
side. The change is the source's finite exponent vector and the
corresponding regenerated powers, masks, coefficient placement, and
evaluation schedule. The Sun-dependent bridge and finalizer are retained
\cite{source40,sun2026}.

The complete formulas and literal division-free circuits are included in
Appendix~\ref{sec:ten-supplement} and the data files. Its ten inputs range independently over
\emph{all} integers. The source-isolation proof is precisely
Section~\ref{sec:sparse}; its carry bounds have the same fixed quadratic
source degree and coefficient bound. Changing the positional weights
does not change the numerical hypotheses of the inherited tail.

The exact degree is
\[
 \boxed{\deg F_{10,W}=297088W+147144.}                \tag{20}
\]
For $W=799$, this recovers the pinned baseline value 237,520,456;
for $W=239$ and 220 it gives 71,151,176 and 65,506,504. The supplied
degree verifier follows the actual fixed-coefficient circuit and checks
the required nonzero leaders. The unrestricted-integer claim retains
the same reliance on Sun's 2026 preprint as the existing ten-integer
construction. The nine-variable positive construction is independent of
that preprint.

\section{Proof dependencies and formalization status}
\label{sec:boundary}

There are three logically different achievements here.

\begin{enumerate}
\item The new finite-data assertions are completely checkable: source
residual replay, isolation of every used monomial, coefficient binding,
gate liveness, operation counts, and degree certificates.
\item The new ordinary proof establishes the sparse extraction, scaled
base bounds, carry coding, and transfer to the fixed source system.
Its unbounded steps are mathematical proofs, not extrapolations from
testing small possible Diophantine solutions.
\item The interpretation of that fixed source as ZFC refutability is
inherited from package 31's compiler proof and the calibrated K215
endpoint. These are cited mathematical dependencies, not consequences
of finite circuit tests.
\end{enumerate}

The active research status at the pinned commit explicitly distinguishes
an unconditional Lean theorem about a classically selected nine-variable
polynomial from the explicit native-window circuits, and lists the
explicit compiler formalization and object-language ZFC derivation as
open tasks \cite{status}. The polynomial constructed here is specified
by computable coefficient data; it is not that classically selected
polynomial. Likewise, a local Lean proof of equality of two zero sets
does not by itself identify either set with ZFC refutations.

Our ordinary ZFC-provability claim uses the imported compiler theorem
in exactly the same way as its source construction. A reader requiring
an end-to-end machine-checked certificate should regard that dependency
as a remaining verification obligation. We neither assume
$\Con$ nor claim to decide it. No inference to an internal ZFC proof is
made merely from the soundness of Lean as a metatheory.

\section{Reproduction, audit design, and further reductions}

The archive is self-contained for its numerical and algebraic checks.
Python's standard library suffices. The release driver checks the
parent data and the materialized ports, rebuilds the selected circuits,
and validates both weight choices. The source verification covers every
used monomial against every possible quadratic expansion monomial, and
also replays all 136 row terms. The degree verification uses an explicit
cancellation identity; it is not an interpolation at a handful of points.
Proof supplements and the immutable source ledger identify which
external theorems are imported.

The polynomial can be evaluated modulo any integer directly from its
JSON register list. Exact evaluation at even small integer assignments
can create impractically large integers, so modular formula checks are
the useful reproducible transcription tests. They do not search for
solutions or establish consistency.

\paragraph{Why this change has room to generalize.}
For a fixed degree-$d$ source, one need only isolate source support
monomials from the degree-$d$ expansion monomials that can contaminate
their coefficient targets. A full $B_d$ code need not be optimal for
that task. Here this distinction is enough to remove most of the
irrelevant pair-separation constraints. The method is not specific to
K215; the source's sparse support is the optimization input.

\paragraph{Most useful next research directions.}
The next finite optimization is to search the actual isolation
constraints with exact solvers or improved local search and to optimize
operation count together with maximum weight. Merely reducing $W$ can
worsen the addition-chain cost, as the two delivered vectors show.
A second possibility is to move the 22 target positions independently
instead of spacing them uniformly by $2W+1$; each proposed packing must
still be checked against cross-row contamination and coefficient
overlap. A third is to change the quadratic source coordinates so that
its used support is easier to isolate. These are further research
questions, not improvements asserted by this release.

The most consequential verification task is separate: formalize the
explicit native compiler and connect its literal coefficient payloads
to the already checked ZFC refutation endpoint. The smaller source
encoding here can then be added as a compact arithmetic theorem with
a finite isolation certificate.

\appendix
\numberwithin{equation}{section}
\section{Complete ten-integer definition and transfer proof}
\label{sec:ten-supplement}

This supplement specifies the second delivered polynomial without any
implicit auxiliary unknowns.  Its ten independent unknowns are
\[
 (f,g,h,k,l,w,x,y,z,m)\in\Z^{10}.
\]
Every other symbol in this section abbreviates an integer polynomial in
these ten variables.  Symbols used for the nine-variable construction
are reset locally.  The parent polynomials $p_j$, their fixed weighted
coefficients $c_{j,\alpha}$, and either one of the two weight vectors
are exactly those in Sections~\ref{sec:parent}--\ref{sec:sparse}.
In particular,
\[
 W=\max_iw_i,\quad h_j=j(2W+1)-1,\quad
 H_*=44W+21,\quad K_*=46W+21,
\]
and $\mathcal D$ is the fixed univariate polynomial in
\eqref{eq:D}.  These are fixed exponents and fixed coefficients,
not further unknowns.

\subsection{The two-input packed source}

Set the fixed integer $\beta_{10}=2^{2140030}$ and define
\begin{align*}
 b&=1+3f,& B_{10}&=\beta_{10}b^2,
 &\ell&=\sum_{i=1}^{31}B_{10}^{w_i},\\
 N_0&=B_{10}^{W+1},&N_1&=4B_{10}^{K_*+1},&N&=N_0N_1,\\
 M_c&=N_0-1-(b-1)\ell,
 &Q_c&=(B_{10}-2)\sum_{j=1}^{22}B_{10}^{h_j},
 &T_c&=M_c+Q_cN_0,\\
 J_c&=2(1+g)^2\mathcal D(B_{10})+
                B_{10}\sum_{q=0}^{K_*}B_{10}^{q},
 &S_c&=g+J_cN_0.
\end{align*}
Finally put
\[
 X=(N-1)\bigl((S_c+T_c+1)N+T_c+1\bigr),\qquad
 Y=N^2,\qquad C_b=bB_{10}^{W}.
\]
Thus $N=4B_{10}^{47W+23}$ and
$C_b=2^{2140030W}b^{2W+1}$.
All sums have fixed finite limits.  The released circuit evaluates them
by shared powers and geometric-sum identities.

\subsection{The retained Lucas and norm registers}

Define, in the displayed order,
\begin{align*}
 L&=lY,&U&=2LX,&V&=4gwY,&W_b&=bw,\\
 A&=U(V+1),&B_1&=2X+1,&C&=B_1+(A-2)h,\\
 P_n&=U^2V,&K&=X+1+k(P_n-2),\\
 S&=2A-5,&T&=3W_bC-2(W_b^2-1),\\
 R_2&=(P_n^2-4)K^2+4.
\end{align*}
The margin and Lucas registers are
\begin{align*}
 \Delta_b&=8C_b^3gK^2-
       g^2\bigl(32(C-KL)^2C_b^3+g^2K^2\bigr),\\
 O&=f^2l^2\Delta_b,&\xi&=zx,\\
 D&=(A^2-4)C^2+4,&E&=C^2D\xi,\\
 F&=4(A^2-4)E^2+1,
 &G^*&=2(A+2)(A-2)^2E^2-1,\\
 \Omega_{10}&=16E,
 &H^*&=C+(B_1-C)F+\Omega_{10}Fy,\\
 I^*&=((G^*)^2-1)(H^*)^2+1.
\end{align*}
The homogeneous quintic merger and the three radicands are
\begin{align*}
 \Sigma&=\xi S^5,&r&=T+Sy,\\
 \Theta_{10}&=(r^3-rS^2-S^3)(r^2+23S^2),\\
 r_0&=b\bigl(1+(3O-4)b^2z^2\bigr),&
 r_1&=R_2,&r_2&=DFI^*,\\
 v_1&=r_0,&v_2&=v_1(P_n^2+K^2+3),\\
 Z_n&=\Sigma m-\Theta_{10},\\
 a_0&=\Sigma^2r_0,&a_1&=\Sigma^2v_1^2r_1,
 &a_2&=\Sigma^2v_2^2r_2.
\end{align*}
Set $V_n=Z_n^2+a_2-a_0-a_1$.  The complete equation is
\begin{equation}
 \boxed{\quad
 F_{10}:=(V_n^2+4Z_n^2a_2-4a_0a_1)^2
              -16V_n^2Z_n^2a_2=0.\quad}
 \label{eq:supplement-F10}
\end{equation}
After the definitions are substituted, this is a single polynomial over
$\Z$ in precisely the ten displayed integer unknowns.  The files
\texttt{sparse220\_ten\_integer.json} and
\texttt{sparse239\_ten\_integer.json} specify the corresponding expanded
binary arithmetic schedules and bind every nonliteral leaf to its exact
materialized coefficient.

\subsection{The precise inherited source interface}

Here is the interface of package 40's source theorem used in the
composition~\cite{source40}.  Let $H_0$ bound the ordinary coefficients
of the shifted parent rows.  One may take
\[
 H_0=8(1+\beta+K_0+M_C+M_F)^2,
 \qquad \beta_{10}>4\cdot33^2H_0.
\]
The displayed inequality follows from the fixed port bit lengths; the
exact coefficient payloads also verify it directly.  Then:
\begin{enumerate}[label=\textup{(\roman*)}]
\item If the $p_j$ have a common zero in $\N^{31}$, then for every
sufficiently large positive integer $f$ for which $b=1+3f$ is both a
square and a power of two, there is a $g$ with
\[
 b\leq g<C_b,\qquad Y\mid\binom{2X}{X}.
\]
\item If $f\ne0$, $b$ is a square and a power of two,
$0\leq g<2C_b$, and $Y\mid\binom{2X}{X}$, then the $p_j$ have a
common zero in $\N^{31}$.
\item Already under $f\ne0$, $b$ a square, and $0\leq g<2C_b$,
without the power-of-two assumption, one has
\[
 0\leq S_c,T_c<N,\quad3\mid X,\quad X>3b,\quad
 Y>\max\{b,256\}.
\]
\end{enumerate}

The new support-isolation property proves the exact same coefficient
identities and coefficient bound used in that theorem.  Distinct
encoded positions, positive leading coefficient 2, target spacing,
and nonnegative exponents have all been retained.  The proof of
\textup{(iii)} bounds $J_c$ for every integer in the broad permitted
interval before any digit recovery.  For \textup{(ii)}, the lower mask
recovers $g=\sum_i n_iB_{10}^{w_i}$ with $0\leq n_i<b$; the upper mask
tests all 22 centered coefficients, and the binary carry identity gives
the binomial divisibility.  For \textup{(i)}, choose arbitrarily large
$b=4^q>\max_i n_i$ and $f=(b-1)/3$, then use the same word.  The parent
zero is nonzero because its last row has constant coefficient 1; thus
$g\geq B_{10}\geq b$.  These observations prove that the source
interface transfers to the two supplied vectors.

Sun's retained argument uses this interface, $b=1+3f$, and
$C_b=2^c b^d$ with $c$ even and $d$ odd.  Here
$c=2140030W$ is even and $d=2W+1$ is odd.  Thus the new source has all
the required congruence and growth hypotheses.  The use is of the
parametric proofs of Propositions 6.1--6.2 and Theorem 4.1 of
Sun's 2026 paper, as composed and documented in packages 31 and
40~\cite{sun2026,source31,source40}.

\subsection{Why the unrestricted norm has exactly the required zeros}

The retained bridge has two relevant consequences.  A common parent
zero supplies positive witnesses with $A>2$ even, $B_1>0$ odd,
$b$ square, $b\leq g<C_b$, $C=u_{B_1}(A)$, $R_2$ square,
$S\mid T$, and $O>0$.  Conversely, for the formulas above,
\[
 b\text{ square},\quad O>0,\quad\xi\ne0,\quad
 DFI^*\text{ square},\quad R_2\text{ square},\quad S\mid T
 \tag{S1}
\]
imply a parent zero.  The reverse proof first obtains the source's
broad ranges from $O>0$; it then recovers the Lucas value and the
power-of-two base.  The full interpolation assertion says that, when
$C=u_{B_1}(A)$ under the positive hypotheses, for any positive $d$
one may choose $\xi>0$ divisible by $d$ such that \emph{every}
residue class modulo $\xi$ contains a $y$ making $DFI^*$ a square.

For completeness, the finalizer enforces exactly (S1) as follows.
The polynomial $\pi(t)=(t^3-t-1)(t^2+23)$ has no rational zero and has
a root modulo every positive integer.  Its cubic discriminant is
$-23$: at an odd prime other than 23, either the cubic has a simple
root or its irreducible Frobenius 3-cycle makes $-23$ a square; at 23
the cubic has simple root 3; at 2 one uses $-23\equiv1\pmod8$.
Lifting and CRT finish the local assertion.  As $S=2A-5$ is always
odd and nonzero, $\Theta_{10}\ne0$ always.  Moreover,
$\Sigma\mid\Theta_{10}$ implies $S\mid T$: if
$\nu_p(S)>\nu_p(T)$, the two homogeneous factors have valuations
$3\nu_p(T)$ and $2\nu_p(T)$, contradicting $S^5\mid\Theta_{10}$.
Conversely, from $S\mid T$ a root of $\pi$ modulo $|\xi|$ prescribes
a residue class for $y$ giving $\Sigma\mid\Theta_{10}$.

For $z\ne0$, the factors $b$ and $1+(3O-4)b^2z^2$ are coprime, and
the second cannot vanish.  Therefore $r_0$ square makes $b$ a signed
square; $b\equiv1\pmod3$ excludes the negative sign.  Its second
factor is then a square, which forces $O>0$.  Conversely positive
bridge witnesses have $O\geq g\geq b\geq4$.  Pell's equation for the
positive nonsquare $3O-4\equiv2\pmod3$, after taking a unit power
whose second coordinate is divisible by $b$, supplies $z>0$ making
$r_0$ square.

Finally, for $z\ne0$ one has $|r_0|>1$ and
\[
 \sqrt{|r_1|}\leq P_n^2+K^2+2,\qquad
 |v_1|>\sqrt{|r_0|},\qquad
 |v_2|>\sqrt{|r_0|}+|v_1|\sqrt{|r_1|}.
\]
These inequalities ensure that a rational signed sum
$\pm\sqrt{r_0}\pm v_1\sqrt{r_1}\pm v_2\sqrt{r_2}$ requires all
three radicands to be integer squares.  Indeed, a multiquadratic-field
automorphism changing a root changes the highest-index changed term
by more than the combined absolute values of all earlier changes.
The argument also covers negative radicands and dependent square
classes.

Equation~\eqref{eq:supplement-F10} is identically the product of its
eight signed-root factors.  If $\Sigma=0$, it equals
$\Theta_{10}^8\ne0$; hence any zero has $\xi=zx\ne0$.  Dividing a
vanishing factor by $\Sigma$ and using separation forces the three
squares and then $\Sigma\mid\Theta_{10}$.  The preceding mergers
give (S1), so the reverse bridge supplies a parent zero.  In the
other direction, choose the positive bridge witnesses, then the Pell
value $z>0$, then an interpolation modulus $\xi$ divisible by $z$,
then a $y$ in the quintic-prescribed residue class.  Put $x=\xi/z$.
All radicands are squares and $\Sigma\mid\Theta_{10}$, so an integer
$m$ makes one norm factor vanish.  This proves
\[
 (\exists\mathbf t\in\Z^{10})\ F_{10}(\mathbf t)=0
 \quad\Longleftrightarrow\quad
 (\exists\boldsymbol n\in\N^{31})\ \bigwedge_j p_j(\boldsymbol n)=0.
\]
Together with the inherited fixed-compiler theorem, classical negation
gives the asserted equivalence with $\Con$ in ZFC.  The argument
retains the cited ordinary proof dependencies and makes no additional
claim of end-to-end proof-assistant certification.

\subsection{Exact degree and supplied evaluators}

Writing $p=47W+23$, the source degrees are
$d_X=6p$, $d_Y=4p$, and $d_C=2W+1$.  The retained tail gives
\begin{align*}
 \deg Z_n&=5d_X+10d_Y+30,\\
 \deg a_0&=14d_X+28d_Y+3d_C+57,\\
 \deg a_1&=26d_X+48d_Y+6d_C+98,\\
 \deg a_2&=114d_X+224d_Y+6d_C+440.
\end{align*}
The last strictly dominates the degrees of $a_0,a_1,Z_n^2$.
Thus the top homogeneous part of $F_{10}$ is the fourth power of the
top homogeneous part of $a_2$, and
\[
 \deg F_{10}=297088W+147144.
\]
The $W=220$ evaluator uses 675 operations and has exact degree
65,506,504; the $W=239$ evaluator uses 662 operations and has exact
degree 71,151,176.  Independent direct leading-coefficient propagation
on these two exported DAGs gives nonzero residues 489390939 and
862326449 modulo 1,000,000,007, respectively, certifying attainment of
the propagated upper bounds.

\clearpage
\section{A secondary improvement: minimum norms at degree 44}
\label{app:minima}

The optional \texttt{minima/} supplement reduces the historical
15-positive-variable, degree-44 sum of squares from 78 binary operations
to 76, or to 75 using its inherited signed index/transport recovery.
The source is package 12's \texttt{ordinary\_gap\_system61.json}, copied
byte for byte from commit \texttt{5b0c9ae9fdae83db563f9d2b61ea0fb1f4a801e6}.
The supplement's \texttt{SOURCE\_PINS.json} records full paths and blob
identities. Both new polynomials bind the complete package-31 coefficient
tuple through local binary payloads.

Take fifteen positive variables
\[
 J,F,\alpha,z,f,h,i,T,s,w,\tau,\eta,\zeta,e,\gamma.
\]
Use the fixed ports $(\beta,K_0,M_C,M_F)$ of Section~\ref{sec:parent}.
Define polynomial abbreviations
\begin{gather*}
 q=\beta J+1,\quad C=q-F-\alpha,\quad X=wq,\quad Y=sq^3,\quad E=XY,\\
 k=\eta+\zeta,\quad c=kY+\eta,\quad a=E+Y,\quad H=4a+3,\quad
 \Delta=a^2+H,\\
 D=X+ac+\gamma H,\qquad
 R=((q-1)(q-F)+\alpha)(q^2-1)+(M_C+qM_F)J,\\
 V=c(Tf-1)-Rf^2,\quad U=\Delta(f^2-1),\quad I=ic^2,\quad W_s=I^2-U,\\
 N_0=\tau^2-XY^2(XY^2+1)k^2,\quad
 N_1=D^2-\Delta c^2,\\
 N_a=U\bigl(V^2-(V+e)^2\bigr)+(V+e)^2,\qquad N_k=k-hE-R.
\end{gather*}
For original transport coordinate $z$ put
$N_t=(K_0+w)C+(q-F)-z(q-1)$. The cheaper equation is
\[
 \boxed{P_{75}=N_0^2+N_1^2+N_a^2+W_s^2+(N_kN_t-1)^2-3=0.}
\]
For the shifted coordinate $z=z_{\rm old}-1$, put
$R_t=(K_0+w)C-F-z(q-1)$. The cleaner local alternative is
\[
 \boxed{P_{76}=N_0^2+N_1^2+N_a^2+W_s^2+(N_k-1)^2+R_t^2-3=0.}
\]

\paragraph{Why subtracting three is sound.}
On every positive assignment,
\[
 N_0^2\geq1,\qquad N_1^2\geq1,\qquad N_a^2+W_s^2\geq1.
\]
Their equality cases are respectively $N_0=1$, $N_1=1$, and
$(N_a,W_s)=(1,0)$.
For the first norm, $n_0=XY^2\geq2$ makes $n_0(n_0+1)$ nonsquare,
so the norm cannot vanish. A norm $-1$ solution
$\xi^2-n_0(n_0+1)\nu^2=-1$ would descend under multiplication by
$2n_0+1-2\sqrt{n_0(n_0+1)}$: its new second coordinate is
$(2n_0+1)\nu-2\xi$, strictly between 0 and $\nu$.
For the main norm, $\Delta=(a+2)^2-1$ is nonsquare; its residue 0 or 3
modulo four also excludes norm $-1$.

If $U>0$, then $U\geq2$. A vanishing auxiliary norm would give
$UV^2=(U-1)(V+e)^2$; since $e>0$, the two roots cannot both vanish.
It would follow that $U(U-1)$ is a square, contrary to its position
between consecutive squares. Thus $N_a\ne0$. If $U=0$, then
$W_s=I^2\geq81$, so the block is larger than 1.
Equality in the block therefore forces $W_s=0$ and $N_a=\pm1$.
Now $U=I^2$, and reduction modulo four makes $N_a$ congruent to
$V^2$ or $(V+e)^2$, excluding $-1$.

Consequently $P_{76}=0$ is tuplewise equivalent to the six historical
comparisons $N_0=N_1=N_a=N_k=1$, $W_s=R_t=0$.
For $P_{75}$ the same argument gives $N_0=N_1=N_a=1$, $W_s=0$,
and $N_kN_t=1$. These give a zero of the ordinary product with
$y=V+e$, which may be signed. The current package-31 theorem in
\texttt{upstream/ORDINARY\_BLANK\_PROOF.md} applies directly on its
nine-positive-coordinate domain and recovers $N_k=N_t=1$
for the authentic native tuple. Package 12's canonical gap theorem
supplies $e>0$ for completeness. This is the ordinary strong equation
$I^2=\Delta(f^2-1)$; it does not assert coordinatewise equivalence
to the normalized strong equation. The semantic compiler and ZFC
interpretation remain inherited dependencies.

\paragraph{Exact finite certificate.}
The 76-operation schedule uses 40 multiplications, 21 additions and
15 subtractions; the 75-operation schedule uses 40, 20 and 15.
Every gate and every supplied input is live. The cancellation
\[
 N_a=V^2-(U-1)e(2V+e)
\]
gives the required degree-22 auxiliary bound. Setting all fifteen
variables equal to $t$, exact symbolic propagation yields degree 44
with leading coefficient
\[
 16\beta^{28}+256\beta^{24}+16\beta^{22}>0.
\]
Thus the total-degree upper bound is attained for the actual coefficients.
The standard-library verifier authenticates all source and payload bytes,
rebuilds both circuits, proves the displayed identity by exact ring
expansion, and checks this leading coefficient. Its calculation uses no
interpolation. This secondary saving is three operations relative to
the historical degree-44 equation; it does not replace the main
nine-variable improvement.


\begin{thebibliography}{9}\small
\bibitem{source31} ProveIt, package 31, \emph{Smaller explicit Diophantine
forms for Con(ZFC)}, coefficient adapter and quadratic system.
Pinned source: \href{\base/con-zfc-diophantine-forms/31-sun10/README.md}{31-sun10}.
\bibitem{source40} ProveIt, package 40, \emph{Quadratic packing},
ten-integer source and circuit, pinned source:
\href{\base/con-zfc-diophantine-forms/40-packing10/README.md}{40-packing10}.
\bibitem{source45} ProveIt, package 45, \emph{Multi-output quadratic
packing}, nine-variable construction and exact degree proof,
\href{\base/con-zfc-diophantine-forms/45-multioutput/nine/MULTITARGET_NINE_PROOF.md}{MULTITARGET\_NINE\_PROOF.md}.
\bibitem{status} ProveIt, \emph{DiophantineConsistency research status},
23 October 2026, pinned \href{\base/STATUS.md}{STATUS.md}.
\bibitem{ratio} ProveIt, \texttt{Jones1982.Nine.ratio\_sound} and
\texttt{ratio\_complete},
\href{https://github.com/VladimirReshetnikov/ProveIt/blob/124c07c0b77a9e20e84676514d8b1ff78536ea3f/Computability/HilbertTenthProblem/Lean/Diophantine/Paper1982/NineRatio.lean}{NineRatio.lean}.
\bibitem{jones} J. P. Jones, \emph{Universal Diophantine equation},
Journal of Symbolic Logic \textbf{47} (1982), 549--571.
\href{https://doi.org/10.2307/2273588}{doi:10.2307/2273588}.
\bibitem{sun2026} Z.-W. Sun, \emph{Ten unknowns for Hilbert's tenth
problem over the integers}, 2026 preprint,
\href{https://arxiv.org/abs/2609.01594v2}{arXiv:2609.01594v2}.
\end{thebibliography}
\end{document}
